\section{Transform: translate, scale, rotate, skew}

Properti \code{transform} mengubah tampilan visual elemen tanpa mempengaruhi layout---elemen yang di-transform tetap menempati ruang aslinya di normal flow. Ini menjadikan transform ideal untuk animasi dan efek hover karena tidak menyebabkan reflow pada elemen lain. CSS menyediakan empat fungsi transform utama yang dapat digabungkan dalam satu deklarasi \cite{mdnCss}.

\begin{table}[htbp]
\centering
\begin{tabular}{|P{3cm}|P{4cm}|P{5cm}|}
\hline
\textbf{Fungsi} & \textbf{Efek} & \textbf{Contoh} \\
\hline
\code{translate(x, y)} & Geser horizontal/vertikal & \code{translate(-50\%, -50\%)} \\
\hline
\code{scale(x, y)} & Perbesar/perkecil & \code{scale(1.1)} \\
\hline
\code{rotate(angle)} & Putar & \code{rotate(45deg)} \\
\hline
\code{skew(x, y)} & Miringkan & \code{skewX(5deg)} \\
\hline
\end{tabular}
\caption{Empat fungsi transform CSS dan contoh penggunaannya}
\end{table}

Beberapa fungsi transform dapat digabung: \code{transform: translateY(-4px) scale(1.02);} menggeser ke atas sekaligus memperbesar. Urutan penulisan penting karena transform diaplikasikan dari kanan ke kiri. Properti \code{transform-origin} mengatur titik acuan transformasi (default: \code{center}); mengubahnya ke \code{top left} mengubah perilaku rotasi dan skala \cite{webdevCss}.

\begin{konsep}
\textbf{Centering dengan Transform.} Teknik centering absolut paling andal menggunakan transform: \code{position: absolute; top: 50\%; left: 50\%; transform: translate(-50\%, -50\%);} --- ini menempatkan elemen tepat di tengah kontainernya tanpa mengetahui ukuran elemen. Meskipun Flexbox centering lebih modern, teknik ini tetap berguna untuk kasus positioning tertentu \cite{cssTricks}.
\end{konsep}

